\documentclass[10pt,a4paper]{amsart}
\usepackage[margin=19mm]{geometry}
\usepackage{amsmath,amssymb,mathtools}
\usepackage[scr=rsfs,cal=euler]{mathalfa}
\usepackage{xfrac}
\usepackage[unicode,hidelinks,bookmarks=false]{hyperref}
\newcommand{\ie}{i.e.\ }
\newcommand{\cf}{cf.\ }
\newcommand{\resp}{respectively\ }
\newcommand{\Subsection}{\subsection}
\providecommand{\nobreakdash}{\nobreak\hbox{-}}
\setlength{\emergencystretch}{2em}
\multlinegap0pt
\usepackage{bm}
\usepackage{bbm}
\usepackage{dsfont}
\usepackage{xspace}
\usepackage{cancel}
\usepackage{mathtools}
\usepackage{amsthm}
\makeatletter
\@ifundefined{theorem}{\newtheorem{theorem}{Theorem}}{}
\@ifundefined{lemma}{\newtheorem{lemma}[theorem]{Lemma}}{}
\makeatother
\newcommand\diag{\operatorname{diag}}
\newcommand{\aspace}{H_{\mathrm{as}}}
\newcommand{\complexes}{\mathbb{C}}
\title[New constructive counterexamples to additivity of minimum ou]{New constructive counterexamples to additivity of minimum output R\'enyi p-entropy of quantum channels}
\author{Krzysztof Szczygielski, Micha{\l} Studzi\'nski}
\date{Source: arXiv:2301.07428v3 (CC BY 4.0)}
\begin{document}
\maketitle

\noindent\emph{Source and licence.} New constructive counterexamples to additivity of minimum output R\'enyi p-entropy of quantum channels. Version arXiv:2301.07428v3 (CC BY 4.0), distributed under the Creative Commons Attribution 4.0 International licence, as recorded in the arXiv licence metadata for this exact version and shown on the abstract page https://arxiv.org/abs/2301.07428v3. Full rights notice: licenses/arxiv-2301.07428-CC-BY-4.0.txt.\par\smallskip

\noindent\textit{Editorial scope.} Counted unit: the Lemma whose source label is \texttt{lemma:ExistenceOfBellStates} in the article (source0.tex lines 416--422, source label \texttt{lemma:ExistenceOfBellStates}), delivered together with the article's own complete original proof of it (source0.tex lines 424--499, one \texttt{proof} environment of 76 lines). No mathematics is written, repaired or re-derived: statement and proof are the article's own text line for line. Required context retained uncounted: eq:BellState (lines 411--413). Every definition, display and label of those lines is retained unchanged. Omitted from this package and not redistributed: 1--410, 414--415, 423--423, 500--1007. Nothing inside a retained excerpt is omitted or altered: every retained body is delivered exactly as the article's lines are written, so no omission receipt of any kind accompanies this package. Citations: the retained excerpts carry no citation command at all, so no bibliography entry of the article is transcribed and no reference is redistributed. Reference adaptation: none was needed and none was made; every reference inside the delivered unit resolves inside it, so no unresolved reference or citation is produced. Printed slips kept literal: no printed word, symbol, hypothesis or number of the counted statement or of its proof was altered. Layout and class: the article's own class is \texttt{amsart} with packages including natbib, inputenc, graphicx, float, tikz, pgfplots, amsmath, amsfonts, amssymb, amsthm, mathrsfs, hyperref; the wrapper is the lane's pinned \texttt{amsart} editorial wrapper at 10pt on A4, which restates the theorem-like environments the retained text needs and adds the article's own macro definitions that the retained text needs (\texttt{\textbackslash aspace}, \texttt{\textbackslash complexes}, \texttt{\textbackslash diag}), with extra wrapper packages (bm, bbm, dsfont, xspace, cancel, mathtools). The delivered numbering is the unit's own. Assets: no retained span contains a figure environment, a graphics-inclusion command, a PDF-page inclusion, a table or a TikZ picture; no image file is copied into this package, the package declares no asset dependency and no image file is redistributed with it. Licence: version arXiv:2301.07428v3 of this article is distributed under the Creative Commons Attribution 4.0 International licence, as recorded in the arXiv licence metadata for this exact version and shown on the abstract page https://arxiv.org/abs/2301.07428v3. A full rights notice is delivered at licenses/arxiv-2301.07428-CC-BY-4.0.txt. Characters: every retained excerpt is pure ASCII, so the delivered engine, XeLaTeX with its default Latin Modern text font, reports no missing character.
\begin{equation}\label{eq:BellState}
    \psi_{r,s} = \frac{1}{\sqrt{d}} \sum_{j=0}^{d-1}\omega^{rj} h_j \otimes f_{(j+s)\,\mathrm{mod}\,d},
\end{equation}

\begin{lemma}\label{lemma:ExistenceOfBellStates}
For each dimension $d \geqslant 2$ there exist 
\begin{equation}\label{eq:sOfDBellStates}
    s(d) = \frac{d}{2} \left[2 + (d+1) \, \mathrm{mod} \, 2\right]
\end{equation}
generalized Bell states $\psi_{r,s}$ in space $\aspace^{\perp}$. Explicitly, we have $s(d) = d$ if $d$ is odd and $s(d) = \frac{3}{2}d$ if $d$ is even.
\end{lemma}

\begin{proof}
Note that every vector $x = \sum_{ij} x_{ij} h_i \otimes f_j \in K\otimes E$ is uniquely identified with matrix $[x_{ij}]$, so it is enough to analyze appropriate matrices instead. Likewise, the full tensor product $K\otimes E$ is isomorphic, as a Hilbert space, with $M_{d}(\complexes)$ equipped with Frobenius (Hilbert-Schmidt) inner product. Then, it is easy to notice that space $\aspace$ is identified with a subspace of all \emph{antisymmetric matrices}, i.e.~matrices $m \in M_{d}(\complexes)$ satisfying $m^T = -m$, and its orthogonal complement $\aspace^{\perp}$ is isomorphic to subspace of all \emph{symmetric matrices}, i.e.~such that $m^T = m$. Then, question of existence of generalized Bell states in $\aspace^{\perp}$ can be rephrased as a problem of finding Bell states described by symmetric matrices; let us call such states simply \emph{symmetric} for brevity. Fix $d \geqslant 2$ and $r \in \{0,\,  \ldots \, , \, d-1\}$. One checks that vector $\psi_{r,0}$ is \emph{diagonal}, i.e.~is  represented by a diagonal matrix
\begin{equation}
    \hat{\psi}_{r,0} = \frac{1}{\sqrt{d}}\diag\left\{1,\, \omega^{r}, \, \omega^{2r}, \, \ldots \, , \, \omega^{(d-1)r}\right\},
\end{equation}
which is clearly symmetric, so we know that there exists at least one Bell state in $H_{\mathrm{as}}^{\perp}$ for each $r$. Prescription \eqref{eq:BellState} allows for easy generation of all vectors $\psi_{r,s}$ for $s > 0$. Below we present a short textual description of this procedure, accompanied by an explicit example in case $d=6$ for reader's convenience. Having matrix $\hat{\psi}_{r,s}$, matrix $\hat{\psi}_{r,s+1}$ is built by cyclically pushing all columns of $\hat{\psi}_{r,s}$ one position to the right (with the last column of $\hat{\psi}_{r,s}$ becoming the first one in $\hat{\psi}_{r,s+1}$). In result, matrix $\hat{\psi}_{r,1}$ will have only one non-zero element $\frac{1}{\sqrt{d}}\omega^{(d-1)r}$ in  position $(d,1)$; likewise, matrix $\hat{\psi}_{r,2}$ will contain $\frac{1}{\sqrt{d}}\omega^{(d-1)r}$ in position $(d,2)$ and $\frac{1}{\sqrt{d}}\omega^{(d-2)r}$ in position $(d-1,1)$ and so forth. In $s$-th step of this process, all columns of initial matrix $\hat{\psi}_{r,0}$ are pushed by $s$ positions to the right and the only non-zero element of the first column of resulting matrix $\hat{\psi}_{r,s}$ lays in $(d-s+1)$-th row. Now, if $\hat{\psi}_{r,s}$ is expected to be non-diagonal and symmetric, number of leading zeros in its first column must equate number of leading zeros in its first row, so we necessarily need $s = d-s$ or $s = \frac{d}{2}$. This is however possible only for even $d$, so we conclude that if dimension $d$ is odd, we only have $s(d) = d$ Bell states in $\aspace^\perp$, precisely \emph{diagonal} ones, $\psi_{0,0}$, $\psi_{1,0}$, \ldots , $\psi_{d-1,0}$. Assume then $d$ is even. In such case, matrix $\hat{\psi}_{r,0}$ may be put in a block form,
\begin{equation}\label{eq:psir0}
    \hat{\psi}_{r,0} = \frac{1}{\sqrt{d}} \left[ \begin{array}{c|c} \mathbf{M}_{11} & \mathbf{0} \\ \hline \mathbf{0} & \mathbf{M}_{22} \end{array} \right],
\end{equation}
where
\begin{align}
    \mathbf{M}_{11} &= \diag\left\{1,\, \omega^{r}, \, \ldots \, , \, \omega^{\left( \frac{d}{2}-1\right) r}\right\},\\
    \mathbf{M}_{22} &= \diag\left\{\omega^{\frac{d}{2}r},\, \omega^{\left( \frac{d}{2}+1\right) r}, \, \ldots \, , \, \omega^{(d-1) r}\right\}.\nonumber
\end{align}
The only candidate, for given $r$, for a symmetric and non-diagonal Bell state is the one for $s = \frac{d}{2}$, so we must have $\hat{\psi}_{r,d/2} = \hat{\psi}_{r,d/2}^{T}$. However, as $\hat{\psi}_{r,d/2}$ was constructed by shifting columns to the right in cyclic manner, one checks that $\hat{\psi}_{r,d/2}$ is simply the block matrix \eqref{eq:psir0} with switched columns,
\begin{equation}
    \hat{\psi}_{r,d/2} = \frac{1}{\sqrt{d}} \left[ \begin{array}{c|c} \mathbf{0} & \mathbf{M}_{11} \\ \hline \mathbf{M}_{22} & \mathbf{0} \end{array} \right].
\end{equation}
The symmetry condition yields then $\mathbf{M}_{11} = \mathbf{M}_{22}$, or, explicitly,
\begin{align}\label{eq:SymmRelation}
    \omega^{rk} = \omega^{r\left(\frac{d}{2} + k\right)} \quad &\text{for all } k \in \{0,\, \ldots \, , \, d-1\},\\
    &\text{or} \quad \omega^{\frac{rd}{2}} = e^{i\pi r} = 1.\nonumber
\end{align}
However, satisfying $e^{i\pi r} = 1$ as appeared in \eqref{eq:SymmRelation} is possible only for \emph{even} $r$. In consequence, a total number of non-diagonal Bell states in $\aspace^{\perp}$ for even values of $d$ equates a number of all even naturals $r \in \{0, \, \ldots \, , \, d-1\}$ which is $\frac{d}{2}$. Together with $d$ diagonal vectors $\psi_{r,0}$ we have $s(d) = \frac{3}{2}d$ symmetric Bell states. Summarizing,
\begin{equation}
    s(d) = \begin{cases}
        d, & d \text{ odd,} \\
        \frac{3}{2}d, & d \text{ even,}
        \end{cases}
\end{equation}
which can be further checked to be equivalent to claimed formula \eqref{eq:sOfDBellStates}. As an example, below we demonstrate the aforementioned procedure in case $d=6$ (with dots representing zeros):
\small
\begin{align}
\hat{\psi}_{r,0} &= \frac{1}{\sqrt{6}}
    \left[
    \begin{array}{ccc|ccc}
        1 & \cdot    & \cdot & \cdot & \cdot & \cdot \\
        \cdot   & \omega^r & \cdot & \cdot & \cdot & \cdot \\
        \cdot   & \cdot    & \omega^{2r} & \cdot & \cdot & \cdot \\
        \hline
        \cdot   & \cdot    & \cdot & \omega^{3r} & \cdot & \cdot \\
        \cdot   & \cdot    & \cdot & \cdot & \omega^{4r} & \cdot \\
        \cdot   & \cdot    & \cdot & \cdot & \cdot & \omega^{5r}
    \end{array} \right] \\
    \xrightarrow{s+1} \hat{\psi}_{r,1} &= \frac{1}{\sqrt{6}} \left[
    \begin{array}{cccccc}
        \cdot & 1 & \cdot & \cdot & \cdot & \cdot \\
        \cdot & \cdot & \omega^r & \cdot & \cdot & \cdot \\
        \cdot & \cdot & \cdot & \omega^{2r} & \cdot & \cdot \\
        \cdot & \cdot & \cdot & \cdot & \omega^{3r} & \cdot \\
        \cdot & \cdot & \cdot & \cdot & \cdot & \omega^{4r} \\
        \omega^{5r} & \cdot & \cdot & \cdot & \cdot & \cdot
    \end{array} \right] \nonumber \\
    \xrightarrow{s+1}\hat{\psi}_{r,2} &= \frac{1}{\sqrt{6}} \left[
    \begin{array}{cccccc}
        \cdot & \cdot & 1 & \cdot & \cdot & \cdot \\
        \cdot & \cdot & \cdot & \omega^r & \cdot & \cdot \\
        \cdot & \cdot & \cdot & \cdot & \omega^{2r} & \cdot \\
        \cdot & \cdot & \cdot & \cdot & \cdot & \omega^{3r} \\
        \omega^{4r} & \cdot & \cdot & \cdot & \cdot & \cdot \\
        \cdot & \omega^{5r} & \cdot & \cdot & \cdot & \cdot
    \end{array} \right] \nonumber \\
    \xrightarrow{s+1}\hat{\psi}_{r,3} &= \frac{1}{\sqrt{6}} \left[
    \begin{array}{ccc|ccc}
        \cdot & \cdot & \cdot & 1 & \cdot & \cdot \\
        \cdot & \cdot & \cdot & \cdot & \omega^r & \cdot \\
        \cdot & \cdot & \cdot & \cdot & \cdot & \omega^{2r} \\
        \hline
        \omega^{3r} & \cdot & \cdot & \cdot & \cdot & \cdot \\
        \cdot & \omega^{4r} & \cdot & \cdot & \cdot & \cdot \\
        \cdot & \cdot & \omega^{5r} & \cdot & \cdot & \cdot
    \end{array} \right]\nonumber
\end{align}
\normalsize
Then, simple analysis shows that the relation \eqref{eq:SymmRelation} for such elements must be met. 
\end{proof}

\end{document}
